\documentclass[12pt,a4paper]{article}
\usepackage[utf8]{inputenc}
\usepackage[english,russian]{babel}
\usepackage{amsmath,amsfonts,amssymb}
\usepackage{geometry}
\usepackage{indentfirst}
\usepackage{hyperref}

\geometry{top=25mm,bottom=25mm,left=30mm,right=15mm}

\title{Решение систем уравнений для нахождения центра и радиуса окружности}
\author{}
\date{\today}

\begin{document}

\maketitle

\section{Задача 1: Нахождение центра окружности заданного радиуса через две точки}

Дана система уравнений относительно центра $(x_0, y_0)$ окружности заданного радиуса $R$, проходящей через две заданные точки $(x_1, y_1)$ и $(x_2, y_2)$:
\begin{equation}
(x_0 - x_1)^2 + (y_0 - y_1)^2 = R^2
\end{equation}
\begin{equation}
(x_0 - x_2)^2 + (y_0 - y_2)^2 = R^2
\end{equation}

\subsection*{Решение}

Геометрически центр окружности $(x_0, y_0)$ равноудален от точек $(x_1, y_1)$ и $(x_2, y_2)$, следовательно, он лежит на серединном перпендикуляре к отрезку, соединяющему эти точки.

\paragraph{Шаг 1: Линеаризация.}
Вычтем из уравнения (1) уравнение (2), чтобы избавиться от квадратичных членов $x_0^2$ и $y_0^2$:
\[
((x_0 - x_1)^2 - (x_0 - x_2)^2) + ((y_0 - y_1)^2 - (y_0 - y_2)^2) = 0
\]
Раскрывая скобки по формуле разности квадратов, получаем уравнение прямой:
\[
2x_0(x_2 - x_1) + x_1^2 - x_2^2 + 2y_0(y_2 - y_1) + y_1^2 - y_2^2 = 0
\]
Разделим на 2 и введем обозначения:
\begin{itemize}
    \item $\Delta x = x_2 - x_1$
    \item $\Delta y = y_2 - y_1$
    \item $x_m = \frac{x_1 + x_2}{2}$ (координаты середины отрезка)
    \item $y_m = \frac{y_1 + y_2}{2}$
\end{itemize}

Тогда уравнение серединного перпендикуляра примет вид:
\[
\Delta x \cdot x_0 + \Delta y \cdot y_0 = \Delta x \cdot x_m + \Delta y \cdot y_m
\]

\paragraph{Шаг 2: Векторный подход.}
Пусть $d = \sqrt{\Delta x^2 + \Delta y^2}$ — расстояние между заданными точками.
\begin{itemize}
    \item Если $d > 2R$, решений нет (точки находятся слишком далеко).
    \item Если $d = 0$, точки совпадают (при $R > 0$ решений бесконечно много).
\end{itemize}

Если решение существует ($d \le 2R$), то центр окружности смещен относительно середины отрезка $(x_m, y_m)$ вдоль перпендикулярного вектора на расстояние $h$:
\[
h = \sqrt{R^2 - \left(\frac{d}{2}\right)^2}
\]

Единичный вектор перпендикуляра к отрезку равен:
\[
\vec{n} = \left(-\frac{\Delta y}{d}, \frac{\Delta x}{d}\right)
\]

\paragraph{Шаг 3: Финальные формулы.}
Смещаясь от середины отрезка в обе стороны на расстояние $h$, получаем два возможных центра:
\begin{equation}
x_0 = \frac{x_1 + x_2}{2} \mp \frac{\Delta y}{d} \sqrt{R^2 - \left(\frac{d}{2}\right)^2}
\end{equation}
\begin{equation}
y_0 = \frac{y_1 + y_2}{2} \pm \frac{\Delta x}{d} \sqrt{R^2 - \left(\frac{d}{2}\right)^2}
\end{equation}
Знаки выбираются попарно: верхний знак для одного центра, нижний — для другого.

\newpage

\section{Задача 2: Нахождение центра и радиуса окружности по трем точкам}

Дана система уравнений относительно неизвестных $(x_0, y_0)$ и $R$ для окружности, проходящей через три точки $(x_1, y_1)$, $(x_2, y_2)$ и $(x_3, y_3)$:
\begin{equation}
(x_0 - x_1)^2 + (y_0 - y_1)^2 = R^2
\end{equation}
\begin{equation}
(x_0 - x_2)^2 + (y_0 - y_2)^2 = R^2
\end{equation}
\begin{equation}
(x_0 - x_3)^2 + (y_0 - y_3)^2 = R^2
\end{equation}

\subsection*{Решение}

Геометрически задача эквивалентна поиску центра описанной окружности треугольника.

\paragraph{Шаг 1: Исключение квадратов и радиуса.}
Вычтем из уравнения (5) уравнение (6):
\[
2(x_2 - x_1)x_0 + 2(y_2 - y_1)y_0 = (x_2^2 - x_1^2) + (y_2^2 - y_1^2)
\]
Вычтем из уравнения (5) уравнение (7):
\[
2(x_3 - x_1)x_0 + 2(y_3 - y_1)y_0 = (x_3^2 - x_1^2) + (y_3^2 - y_1^2)
\]

Получена система линейных уравнений вида:
\[
\begin{cases}
A_1 x_0 + B_1 y_0 = C_1 \\
A_2 x_0 + B_2 y_0 = C_2
\end{cases}
\]
где:
\begin{align*}
A_1 &= 2(x_2 - x_1), & B_1 &= 2(y_2 - y_1), & C_1 &= (x_2^2 - x_1^2) + (y_2^2 - y_1^2) \\
A_2 &= 2(x_3 - x_1), & B_2 &= 2(y_3 - y_1), & C_2 &= (x_3^2 - x_1^2) + (y_3^2 - y_1^2)
\end{align*}

\paragraph{Шаг 2: Метод Крамера.}
Вычислим определители:
\[
D = A_1 B_2 - B_1 A_2 = 4 \cdot \left((x_2 - x_1)(y_3 - y_1) - (y_2 - y_1)(x_3 - x_1)\right)
\]
Если $D = 0$, точки лежат на одной прямой, окружности не существует. Если $D \neq 0$:
\[
D_x = C_1 B_2 - B_1 C_2, \quad D_y = A_1 C_2 - C_1 A_2
\]
Координаты центра:
\[
x_0 = \frac{D_x}{D}, \quad y_0 = \frac{D_y}{D}
\]

\paragraph{Шаг 3: Явные аналитические формулы и радиус.}
В развернутом виде:
\begin{equation}
x_0 = \frac{(x_1^2 + y_1^2)(y_2 - y_3) + (x_2^2 + y_2^2)(y_3 - y_1) + (x_3^2 + y_3^2)(y_1 - y_2)}{2 \cdot \left(x_1(y_2 - y_3) + x_2(y_3 - y_1) + x_3(y_1 - y_2)\right)}
\end{equation}
\begin{equation}
y_0 = \frac{(x_1^2 + y_1^2)(x_3 - x_2) + (x_2^2 + y_2^2)(x_1 - x_3) + (x_3^2 + y_3^2)(x_2 - x_1)}{2 \cdot \left(x_1(y_2 - y_3) + x_2(y_3 - y_1) + x_3(y_1 - y_2)\right)}
\end{equation}

После нахождения $(x_0, y_0)$ радиус $R$ вычисляется из любого исходного уравнения:
\begin{equation}
R = \sqrt{(x_0 - x_1)^2 + (y_0 - y_1)^2}
\end{equation}

\end{document}
